%!TEX root = ../template.tex
%%%%%%%%%%%%%%%%%%%%%%%%%%%%%%%%%%%%%%%%%%%%%%%%%%%%%%%%%%%%%%%%%%%%
%% chapter-related-work.tex
%% Grande Panorama de Lisboa — AR Companion App
%% Chapter 2: Related Work — SKELETON (~18 pages)
%%%%%%%%%%%%%%%%%%%%%%%%%%%%%%%%%%%%%%%%%%%%%%%%%%%%%%%%%%%%%%%%%%%%

\typeout{NT FILE chapter-related-work.tex}%

\chapter{Related Work}
\label{ch:related-work}


% ===================================================================
\section{Digital Heritage and Interactive Museum Experiences}
\label{sec:rw-heritage}
% ===================================================================

Museums have spent roughly four decades renegotiating their relationship
with the visitor, moving from an institution that displays objects toward
one that designs experiences. This section traces that shift, situates
digital and interactive tools within it, and reviews what visitor-behaviour
research tells us about attention, motivation, and fatigue. These three
constraints end up shaping almost every design decision in our own AR
panorama application.

\subsection{Evolution of Museum Interpretation}
\label{subsec:rw-museum-evolution}

The so-called \emph{New Museology}, articulated by Vergo and collaborators
\cite{vergo1989}, reframed the museum's core task. Where earlier practice
treated the object and its label as sufficient, this paradigm argued that
meaning is something visitors construct, not something a label transmits;
the institution's job shifted from safeguarding objects to designing for
that construction of meaning. That is not just a historical footnote. It
is the theoretical justification for everything that follows in this
chapter, since AR-mediated interpretation only makes sense once one
accepts that the visitor's experience, not the artefact alone, is the
unit being designed.

Successive waves of interpretive technology followed this visitor-centred
logic without abandoning it: the audio guide, then the interactive kiosk,
then the dedicated mobile app, and now mobile AR each added interactivity
while generally preserving physical proximity to the artefact. We read
this less as a simple upgrade path than as a recurring negotiation
between depth of interpretation and fidelity to the physical encounter. Recupero
et al.\ \cite{recupero2019} make this tension concrete at Rome's Ara
Pacis Museum, analysing an AR/VR-augmented tour through Activity Theory
and finding that the technology succeeded when it mediated the visitor's
relationship with the artefact, and generated friction when it competed
with it for attention. For a project built around a single architectural
surface, this finding is directly actionable: every interface decision
should be tested against whether it deepens the encounter with the tile
wall or distracts from it.

\subsection{Digital Technologies in Cultural Heritage}
\label{subsec:rw-ch-tech}

Cultural heritage institutions now draw on a fairly wide spectrum of
digital tools, ranging from static aids such as QR codes and audio guides,
through touchscreen kiosks and dedicated apps, to spatially aware formats
such as projection mapping and mobile AR. The axis that matters most for
our purposes is not the technology itself but whether it is
\emph{passive} (fixed content, the same for every visitor) or
\emph{interactive} (responsive to where the visitor stands and what they
are looking at). Mobile AR sits at the interactive end of that spectrum
almost by definition, since it has no function at all without first
detecting the visitor's position relative to the artefact.

Two bibliometric surveys give us a sense of how fast this interactive end
of the spectrum has grown. Boboc et al.\ \cite{boboc2022} analysed 1,201
documents indexed in Scopus and Web of Science between 2012 and 2021 and
identified eight recurring research themes, among them 3D reconstruction,
virtual museums, education, tourism, and gamification. That spread
suggests AR in cultural heritage is no longer a niche experiment but a
multi-pronged research area with its own internal specialisations.
Bekele et al.\ \cite{bekele2018}, in an earlier and more technically
focused survey, organise the field around three recurring design
concerns (tracking and registration, virtual-content modelling, and presentation/interaction design), which maps closely onto the three
technical problems our own panorama application must solve: detecting
150 POIs reliably, modelling their associated content, and presenting that
content without overwhelming the visitor.

Whether this interactivity actually helps visitors is, fortunately, not
just a theoretical claim. Bachiller, Monzo, and Rey \cite{bachiller2023}
compared traditional and AR-augmented visits at a technological heritage
museum and found that more than 90\% of participants who took the
AR-based route reported positive satisfaction, with a statistically
significant association between visit type and both satisfaction and
self-reported learning. Wang and Zhu \cite{wangzhu2022}, reviewing mobile
AR museum applications more broadly, frame the underlying problem these
results are solving: traditional museum displays tend to rely on a single,
fixed interaction style, which raises cognitive load and leaves visitors
with a comparatively poorer experience. Mobile AR addresses this largely
by letting the information layer adapt to the visitor's viewpoint instead
of the reverse. That is exactly the affordance a 23-metre,
150-POI panorama wall needs, since no single fixed sequence of stops could
reasonably serve every visitor's path through it.

\subsection{Visitor Behaviour and Museum Learning Models}
\label{subsec:rw-visitor-models}

Falk and Dierking's Contextual Model of Learning \cite{falkdierking2000}
remains the most widely used framework for explaining \emph{why} two
visitors leave the same exhibit with entirely different experiences. The
model holds that a visit is shaped by three interacting contexts: personal (prior knowledge, motivation, identity), social (companions, shared meaning-making), and physical (the space and the artefact itself). These contexts are continuously constructed by the visitor,
not fixed by the institution. For an AR application, this
implies that content cannot be authored for an abstract "average visitor";
it has to leave room for the personal and social contexts to shape how
each visitor actually uses it.

Falk's later work on identity \cite{falk2009} operationalises the personal
context into five recurring visitor identities (explorer, facilitator, experience seeker, professional/hobbyist, and recharger) each with a
different relationship to depth of content and pacing. Around the same
period, Pekarik, Doering, and Karns \cite{pekarik1999}, working from
interviews across nine Smithsonian museums, proposed a complementary
four-category model of what visitors actually find satisfying: object
experiences (encountering the "real thing"), cognitive experiences
(interpretive depth), introspective experiences (personal reflection),
and social experiences (shared with companions). Taken together, these
two models suggest that a single POI in our panorama should ideally be
able to serve more than one of these motivations: a quick factual layer
for the experience seeker, say, alongside an optional deeper layer for
the professional or the introspective visitor, instead of committing to
one fixed tone for all 150 points.

This flexibility matters because attention in museum settings is a scarce
resource and a fairly fragile one. Serrell's large-scale observational
study across 108 exhibitions \cite{serrell1997} found that visitors
typically spend under twenty minutes in any given exhibition regardless of
its size or topic, and that most visitors do not stop at more than half of
the available elements, which should temper any assumption
that visitors will systematically work through all 150 POIs in sequence.
Bitgood's reviews of what is generally termed \emph{museum fatigue}
\cite{bitgood2009a, bitgood2009b} describe \emph{satiation} as the
dominant mechanism behind this decline: attention and enjoyment both drop
after repeated exposure to similar stimuli, independent of physical
tiredness. Antón, Camarero, and Garrido \cite{anton2018} add a design-relevant
nuance to this picture: visitors who follow a self-regulated, free route
through a museum appear to reach satiation later, and report higher
satisfaction, than those guided along a fixed, ordered path. Table
\ref{tab:visitor-behaviour-design} summarises how these findings translate
into concrete constraints for a 150-POI panoramic AR experience.

\begin{table}[h]
\centering
\caption{Visitor-behaviour findings and their design implications for a
150-POI panorama AR application.}
\label{tab:visitor-behaviour-design}
\small
\begin{tabular}{p{7cm}p{7cm}}
\toprule
\textbf{Finding} & \textbf{Design implication} \\
\midrule
Visitors rarely stop at more than half the available elements; sessions
average under 20 minutes (Serrell \cite{serrell1997}) &
Do not design for full 150-POI coverage; support short, self-contained
sub-sessions \\
\addlinespace
Satiation lowers attention/enjoyment after repeated similar stimuli,
independent of physical fatigue (Bitgood \cite{bitgood2009a,bitgood2009b}) &
Vary POI content format (text, audio, short reconstruction) instead of
repeating one template 150 times \\
\addlinespace
Free, self-regulated routes delay satiation and increase satisfaction
relative to fixed routes (Antón, Camarero \& Garrido \cite{anton2018}) &
Let visitors choose their own path across the wall instead of imposing a
single guided sequence \\
\addlinespace
Visitors pursue distinct motivations (explorer, experience seeker,
recharger, etc.) (Falk \cite{falk2009}) &
Layer content depth per POI so each motivation type can self-select how
much it engages with \\
\bottomrule
\end{tabular}
\end{table}

























% ===================================================================
\section{Augmented Reality in Cultural Heritage}
\label{sec:rw-ar-ch}
% ===================================================================

Augmented reality did not arrive in museums as a single technology so
much as a loose family of techniques unified by one property: the
physical object stays in view while a digital layer is added around it.
This section first fixes what counts as AR and how it is tracked
(\S\ref{subsec:rw-ar-concepts-tracking}), then reviews how museums have
actually deployed it (\S\ref{subsec:rw-ar-museum-apps}). Throughout, we
treat the Grande Panorama de Lisboa as our primary worked example, but we
draw where useful on the Chafariz Velho fountain wall in Pa\c{c}o de
Arcos (a contemporaneous, tile-covered, and geometrically harder heritage surface used extensively during our own development) as a
reminder that the problems addressed in this thesis generalise beyond
one specific panel.

\subsection{AR Concepts, Taxonomy and Tracking Techniques}
\label{subsec:rw-ar-concepts-tracking}

Milgram and Kishino's reality--virtuality continuum \cite{milgram1994}
places augmented reality close to the physical end of a spectrum that
runs from the unmediated real world through to fully synthetic virtual
environments. AR sits near that real end on purpose: the original
artefact remains visible and dominant, and the virtual layer is added,
not substituted, which is the property that makes AR
attractive in a heritage context where the institution's primary duty is
to the physical object. Azuma \cite{azuma1997} narrows the definition
further into three properties a system must satisfy to count as AR at
all: it combines real and virtual content, it interacts in real time,
and it registers virtual objects within three-dimensional physical
space. None of the three is trivial once an irreplaceable heritage
surface is involved: the artefact cannot be altered to assist
tracking, registration has to survive a visitor walking the length of
the wall, and, as discussed below, the geometry of the surface itself
constrains which tracking strategy is even feasible.

Three broad device categories implement these properties in practice.
\textbf{Head-mounted AR} offers the most immersive registration but remains
largely absent from public museum settings; cost, hygiene concerns
around shared headsets, and a degree of social stigma keep it niche even
in recent surveys of the field \cite{wang2025grand}. \textbf{Spatial AR}, typically
projection mapping onto a fa\c{c}ade or gallery wall, produces a shared
experience without requiring visitors to carry any device, but it
depends on a fixed installation and therefore cannot travel with a
temporary exhibition or a budget-constrained museum. \textbf{Mobile AR}, run on
the visitor's own smartphone, sacrifices some of that immersion but asks
nothing of the host institution beyond Wi-Fi: no hardware purchase, no
headset queue, no installation footprint on a protected surface. For an
application intended to generalise across heritage walls of very
different sizes, mobile AR is the only category that scales without
site-by-site capital investment, which is why it is the device class
this thesis adopts.

Within mobile AR, four tracking strategies recur in the literature, and 
each carries a different cost in a heritage setting. \textbf{Marker-based}
tracking (QR codes, ArUco fiducials, printed image targets) gives
the most reliable registration but requires physically attaching
something to the artefact, which conservation protocols on a protected
tile surface will rarely permit. \textbf{Markerless visual-inertial SLAM}, the
approach underlying both ARKit and ARCore, instead fuses camera feature
points with inertial sensor data and needs no physical modification at
all; its weakness is that tracking quality degrades on featureless or
highly repetitive surfaces, a genuine concern for a tile wall whose
pattern repeats every few panels. \textbf{Area scanning or spatial mapping}
(Vuforia Area Targets, Immersal's cloud-hosted VPS) builds a persistent
3D anchor from a full-room scan and survives across sessions.
\textbf{GPS-based geospatial tracking} sidesteps the surface entirely by anchoring
content to Street-View-derived coordinates instead; Bauer, Hagler, and
Kocur's recent outdoor mural application at Mural Harbor in Linz
\cite{bauer2025} demonstrates this convincingly for fa\c{c}ade-scale art,
but the approach simply does not function indoors, ruling it out for a
gallery setting. Geometry adds a further complication that the
literature rarely makes explicit: a flat panel only needs $x,y$ tracking
along its length, a U-shaped wall such as the Grande Panorama
additionally needs $z$ depth as its facing direction changes between
panels, and a free-standing circular wall such as Chafariz Velho needs
yaw rotation on top of that as a visitor walks its full circumference.
A tracking approach validated only on a flat or U-shaped wall does
not automatically generalise to a circular one. Table
\ref{tab:tracking-comparison} summarises these trade-offs.

\begin{table}[h]
\centering
\caption{Mobile AR tracking strategies and their constraints for
heritage-wall applications.}
\label{tab:tracking-comparison}
\small
\begin{tabular}{p{3.4cm}p{2.0cm}p{2cm}p{2.7cm}p{2.9cm}}
\toprule
\textbf{Strategy} & \textbf{Physical modification} & \textbf{Persistence} & \textbf{Setup burden} & \textbf{Curved surfaces} \\
\midrule
\textbf{Marker-based} & Required & High & Low & Yes \\
\textbf{Markerless VI-SLAM} & None & Session-only & Low & Yes, with denser zones \\
\textbf{Area scan / VPS} & None & High & High (controlled scan) & Yes \\
\textbf{GPS / Geospatial} & None & High & Low--Medium & Outdoor only \\
\bottomrule
\end{tabular}
\end{table}

The SDK landscape implementing these strategies has consolidated
noticeably in recent years. ARKit and ARCore remain the two OS-native
runtimes underneath almost everything else, offering image tracking,
plane detection, and six-degrees-of-freedom pose estimation out of the
box \cite{wangzhu2022}. Unity's AR Foundation sits above both as a
cross-platform abstraction layer, letting a single codebase target iOS
and Android without forking the tracking logic, and has become the de
facto standard for cross-platform mobile AR development largely for that
reason \cite{wangzhu2022}. Flutter is a newer entrant to the AR space, and its ecosystem is still maturing; it allows for faster development and easier cross-platform support than Unity, but it does not yet have a first-party AR abstraction layer and relies on community plugins that bridge separately to ARKit and ARCore ~\cite{arflutterplugin2026}.
Vuforia and Immersal occupy a narrower niche:
Vuforia for high-precision image-target recognition and room-scale Area
Targets, Immersal for cloud-hosted or fully
on-device visual positioning~\cite{immersal2026faq}.
We return to this comparison in Chapter~3, where the choice between an image-tiling
approach and an area-scan approach is made concrete for an actual
heritage wall.

\subsection{AR Applications in Museums and Heritage Sites}
\label{subsec:rw-ar-museum-apps}

How much of this technical machinery has actually reached museum floors?
Wang et al.'s 2025 scoping review \cite{wang2025grand} offers the most
current answer available: starting from 5{,}368 candidate articles
indexed across the ACM Digital Library, IEEE Xplore, and Scopus up to
January 2024, the authors narrowed their sample to 177 empirical studies
through a multi-stage screening process. Exhibition and museum-experience
enhancement turned out to be the single largest application area,
accounting for 61 of 195 coded application instances (31.2\%), ahead of
education and learning at 24.1\%; AR and VR devices together were the
most frequently cited device category at 29.1\%, although mobile and PC
platforms remained the broadest deployment channel overall. Publication
volume climbed steadily after 2017 and peaked at roughly 30 articles in
2023 alone, a reasonably strong signal that immersive cultural-heritage
research is still accelerating, not levelling off.

Within museum-specific mobile AR, Wang and Zhu's survey \cite{wangzhu2022}
organises applications into three recurring functions: \textbf{sensory
stimulation}, where AR simply enriches what a visitor already sees; \textbf{guide
design}, where it supports navigation and wayfinding; and what the authors
term a \textbf{heart-flow experience}, combining game mechanics with interactive
narration. Evidence on outcomes generally favours the richer end of that
spectrum. Bachiller, Monzo, and Rey \cite{bachiller2023} compared
AR-augmented and traditional visits at a technological heritage museum
in Spain and found that more than 90\% of AR visitors reported positive
satisfaction, with a statistically significant association between visit
type and both satisfaction and self-reported learning gain. This kind of
direct comparison remains uncommon in this literature, since most studies
evaluate the AR condition in isolation, without a traditional
baseline. Chen and Lai \cite{chenlai2021}, applying the ARCS motivation
framework to a museum AR deployment, similarly found that the
\emph{relevance} of content to a visitor's own interests predicted
learning effectiveness more strongly than novelty alone, which suggests
that adding AR to an exhibit buys comparatively little on its own;
tailoring what that AR actually shows appears to matter more.

A Portuguese precedent is directly instructive here. Gkatsou's 2018
dissertation at the University of Porto \cite{gkatsou2018} developed a
non-linear mobile AR and QR-code experience for teenage visitors at the
Alf\^{a}ndega Museum of Transports and Communications, combining
storytelling and gamification and validating the design through an
expert panel spanning experience design, visitor studies, and gaming.
The non-linear structure was a deliberate response to a constraint our
own project shares: a museum cannot enforce a fixed visiting order in a
real gallery, so content has to remain coherent regardless of which point
a visitor reaches first. That same constraint, multiplied across roughly
150 points of interest on the Grande Panorama, not just a handful of
exhibition stops, is one of the central design problems this thesis
returns to in the following subsection.









\subsection{AR for Large-Scale and Panoramic Heritage Surfaces}
\label{subsec:rw-large-scale}
 
Almost everything surveyed so far assumes a heritage object small enough
to fit inside a single camera frame, or a fixed installation built around
one façade. The Grande Panorama de Lisboa (and, just as importantly, the considerably larger Chafariz Velho fountain wall in Pa\c{c}o de Arcos, which we relied on extensively during development while the Museu Nacional do Azulejo remained closed for renovation) break that
assumption simply through length. Both are continuous, tile-covered,
mid-eighteenth-century surfaces several tens of metres long, and neither
fits the implicit unit of analysis in most prior AR-heritage work: one
object, one fixed viewpoint, one short interaction. This subsection asks
what changes once the surface itself becomes the scale problem.
 
Three technical issues recur once length enters the picture. First, a
single reference image cannot capture a 20--40\,m painted wall at the
resolution feature-matching needs; down-sampling to one image collapses
feature density below a usable threshold, which makes some form of
multi-zone tiling close to unavoidable for a panel at this scale. Second,
inertial drift in markerless VI-SLAM accumulates with horizontal travel
distance, so a tracking solution validated on a small room rarely
transfers cleanly to tens of metres of lateral walking without
re-anchoring. Third, the glazed tile surface common to both the Grande
Panorama and Chafariz Velho introduces an adversarial visual
environment: specular highlights shift with the time of day, and the
painted scene itself is visually dense enough to compete with any
overlay for attention. This is not a minor concern. Scargill et al.
\cite{scargill2024}, in a controlled study of environment-texture effects
on AR task performance, found that visual background complexity
significantly slowed task completion ($p = 0.003$) without any
compensating gain in how stable the virtual content appeared to users,
which suggests that a visually busy heritage surface actively works
against AR legibility, not for it.
 
Two further dimensions of the problem are worth separating from this
technical core. The first is purely geometric: a flat panel only needs
$x,y$ tracking along its length, the Grande Panorama's U-shaped
arrangement additionally needs $z$ depth as the wall's facing direction
changes between panels, and Chafariz Velho's free-standing circular form
needs yaw rotation on top of that to keep content aligned as a visitor
walks its full circumference. This distinction matters for any claim
of generalisability, since a tracking approach proven on a flat or
U-shaped wall does not automatically extend to a circular one. The second
is what we call the multi-scale viewing
problem: within the same visit, the same painted surface can be read from afar, 
e.g. 3\,m away for a full overview of a scene, and nearby the wall from under 1\,m 
for a single building's detail, and feature detection has to survive that transition
without losing lock.
 
Existing AR work on large or panoramic surfaces offers partial answers
at best. Spatial AR (projection mapping at room or façade scale)
is the closest conceptual match; Printezis and Koutsabasis's 2026
critical review of 52 SAR studies in arts and culture \cite{printezis2026}
proposes a clean three-scale taxonomy (object/tabletop, room/gallery,
building/outdoor) and finds that room-scale installations are
particularly well suited to dialogic, chaptered narratives shared by
co-present visitors without any device at all. That same review
classifies "time-based palimpsests" (layering different historical eras over one physical surface) among the recognised narrative
patterns in this space, which is encouraging for any future timeline
feature built on top of a tile wall. The catch is the same one raised in
\S\ref{subsec:rw-ar-concepts-tracking}: SAR needs a fixed projector
installation, which rules it out for a portable, BYOD-first application.
Outdoors, Bauer, Hagler, and Kocur's mural application at Mural Harbor in
Linz \cite{bauer2025} is arguably the closest published precedent in
surface type (building-scale painted art, tracked markerlessly) but
its GPS/Geospatial approach depends on outdoor Street View coverage and
simply has no indoor equivalent. Sauter et al.'s AR-Explorer
\cite{sauter2024} pushes a related idea further, retrieving and
overlaying historical photographs at their original outdoor capture
locations, but again the anchoring is location-based, not
image-based, and the deployment context is outdoor public space, not
an indoor gallery wall. None of these precedents combines a
continuous, indoor, image-anchored heritage surface at panorama scale
with the multi-zone hand-off problem our wall requires; Table
\ref{tab:large-scale-precedents} makes that gap explicit.
 
\begin{table}[h]
\centering
\caption{Closest published precedents for large-scale heritage AR and
their relevance to a continuous indoor tile-panorama application.}
\label{tab:large-scale-precedents}
\small
\begin{tabular}{p{2.3cm}p{2.6cm}p{2cm}p{2.6cm}p{3.2cm}}
\toprule
\textbf{Work} & \textbf{Surface type} & \textbf{Setting} & \textbf{Tracking} & \textbf{Gap vs.\ this thesis} \\
\midrule
Bauer et al.\ \cite{bauer2025} & Building-scale murals & Outdoor & GPS / Geospatial & No indoor equivalent \\
Sauter et al.\ \cite{sauter2024} & Historical photographs at sites & Outdoor & Location-based & Location-, not image-anchored \\
Printezis \& Koutsabasis \cite{printezis2026} & Room/façade installations & Indoor \& outdoor & Fixed projection & Not portable / BYOD \\
Xu et al.\ \cite{xu2026panoramic} & Panoramic cave murals & Indoor (offline 3D) & Photogrammetry from existing images & Not a live AR tracking system \\
\bottomrule
\end{tabular}
\end{table}
 
\subsection{Research Gap: Multi-POI Detection on a Single Panoramic Surface}
\label{subsec:rw-gap}
 
Drawing the previous two subsections together, four distinct gaps
emerge. None of them is specific to the Grande Panorama; all apply to
any sufficiently large, continuous heritage wall, whether flat, U-shaped
like the Panorama, or curved like Chafariz Velho.
 
The first is technical: no published mobile AR setup has been
demonstrated to maintain stable tracking across a single, continuous
indoor wall of this length, flat or curved, without losing alignment or
drifting as the visitor walks its length. Most AR tooling, and most of
the literature reviewed above, targets small objects or room-scale
environments, not a single elongated, reflective, repetitively
patterned surface; this thesis directly compares a feature-matching
zone-tiling approach against a geometric area-scan approach to establish
which copes better with a wall at this scale, and, where Chafariz Velho
data is available, at this geometry.
 
The second concerns historical layering. Existing heritage AR
applications generally commit to one fixed reconstruction of one
historical moment, with no way for a visitor to move between several;
we are not aware of one that ties a multi-era timeline directly to coordinates
on both a physical wall and a modern map of the surrounding city in a
single synchronised interface. Where time permits, this thesis explores
a state-managed timeline (pre-1755, the 1755 earthquake, the Pombaline
reconstruction, and the present day) anchored to the same coordinate
system used for wall tracking, generalisable to any heritage surface with
a comparably layered history.
 
The third is evaluative. Most empirical studies cited in this chapter
evaluate learning and engagement on small exhibits, single artefacts, or
short VR sessions; we are not aware of comparable real-world data for a
large, narrative-dense tile mural specifically. A central empirical
contribution of this thesis is gathering and analysing in-app analytics,
brief exit surveys, and where feasible a short follow-up questionnaire
from real museum visitors, to begin closing that gap for this class of
heritage object, not for any one site.
 
The fourth is personalisation at scale. Section~\ref{subsec:rw-visitor-models}
established that visitors arrive with different motivations and
attention budgets, yet most AR frameworks still offer one fixed content
depth and one fixed navigation structure regardless of who is using them.
With roughly 150 points of interest available on the Grande Panorama
alone, a one-size-fits-all interface risks the satiation effects
discussed earlier well before a visitor reaches even a fraction of them;
this thesis implements a lightweight onboarding step that adapts content
depth and route suggestions to whether the visitor identifies, loosely,
as a casual tourist, a student, or a more specialised researcher.
 
Taken together, these four gaps motivate a single underlying claim: this
thesis develops, to our knowledge, the first mobile AR framework
explicitly designed for a continuous, large-scale heritage wall (generalisable across flat and curved geometries alike) combining
multi-zone tracking, optional historical-timeline layering, and
visitor-adaptive content within one application. The Grande Panorama
remains the motivating case throughout, and Chafariz Velho the harder
geometric test case used during development, but neither is treated as
the limit of what the resulting framework can support.



























% ===================================================================
\section{UX, Engagement and Personalisation in Heritage AR}
\label{sec:rw-ux}
% ===================================================================

Having established what AR can technically do on a heritage wall, this section
turns to whether visitors actually want to use it, and how. As before, the Grande
Panorama remains the motivating case and Chafariz Velho the harder geometric test
case, but the design principles discussed here are argued on purpose at the
level of ``a large heritage wall a visitor walks past,'' not at the level of one
specific museum.

\subsection{User-Centred Design Principles for Heritage AR}
\label{subsec:rw-ucd}

User-centred design, formalised in ISO~9241-210 \cite{iso9241210}, treats interface design as an
iterative cycle (understand context, specify requirements, produce design solutions, evaluate against requirements, repeat), not a single build
followed by a single test. Applied to a heritage AR app, this means involving
curators and visitor-studies staff throughout development, not only at a final
validation stage, and it means Norman's classic usability principles
\cite{norman1988} translate fairly directly into concrete interface rules: tap
targets need visible affordances, every interaction needs immediate visual or audio
feedback, the number of simultaneous on-screen choices needs limiting to avoid
decision paralysis, and any pin or marker placed on the wall needs to map
unambiguously onto the physical building or scene it annotates.

One principle recurs across the heritage AR literature specifically: virtual content
performs better when it is anchored to the physical surface itself instead of
floated as a screen-space overlay. This follows directly from Azuma's registration
requirement (\S\ref{subsec:rw-ar-concepts-tracking}), but it is also an empirical
finding in its own right. Recupero et al.\ \cite{recupero2019}, studying an
AR/VR-augmented tour at Rome's Ara Pacis Museum, found that the technology worked
best when it visibly mediated the visitor's relationship with the artefact and
generated friction whenever it competed with the artefact for attention, which is
the failure mode a screen-floating overlay risks. We treat this as a hard design
constraint: every piece of AR content in this thesis is
anchored to a specific building or scene on the wall, never presented as a free-
floating screen element.

A second, related principle concerns how much content to reveal at once. Wang
et al.\ \cite{wang2025grand} identify information overload (too many simultaneous virtual elements competing for attention) as a recurring failure mode across the
immersive cultural-heritage studies in their review, one that measurably reduces
visitor focus instead of enriching it. A progressive-disclosure structure, where a
short initial layer is visitor-extendable into deeper detail, addresses this directly; whether visitors actually choose to extend
it is itself an interesting empirical question this thesis can speak to. The same
logic applies to how much of a wall a visitor is expected to engage with at all:
§\ref{subsec:rw-visitor-models} already established, via Bitgood's account of museum
satiation \cite{bitgood2009a,bitgood2009b} and Ant\'on, Camarero and Garrido's finding
that self-regulated routes delay satiation \cite{anton2018}, that visitors disengage
well before exhausting a large exhibit, and that letting them choose their own path
helps more than it hurts. A wall with 150 points of interest, or one with several
dozen spread across 30--40\,m, makes this a first-order design constraint. Gkatsou's expert-validated finding that non-linear AR
structures suit museum visitors better than enforced sequences
\cite{gkatsou2018} reinforces the same conclusion from a different angle: a visitor's
route through a physical space cannot realistically be controlled, so the content
structure has to tolerate (and ideally reward) arriving at any point first.

Aesthetic integration deserves separate treatment because heritage AR carries a risk
most AR design guidance does not address: visual dissonance between a modern
interface and a centuries-old surface. He and Wang \cite{hewang2025}, surveying 324
museum AR visitors, found aesthetics to be a significant predictor of perceived
enjoyment ($\beta = 0.234$, $p < 0.01$), and on a surface as visually dense as a
glazed tile mural, this is more than a stylistic nicety. Scargill et al.\
\cite{scargill2024} demonstrate the cost of getting it wrong directly: visually
complex backgrounds significantly slow AR task completion ($p = 0.003$) without any
corresponding improvement in how stable virtual content appears to the user. Since
the tile pattern itself functions as both the tracking reference and the visual
backdrop on a wall like the Grande Panorama or Chafariz Velho, it cannot simply be
simplified away; the available lever is strong contrast and a spare
UI vocabulary that does not compete with the painted scene for attention. This
connects to a broader principle Printezis and Koutsabasis articulate in their review
of spatial AR storytelling \cite{printezis2026}: use the smallest technological
footprint the cultural purpose requires. If static text and a single image
adequately tell a point of interest's story, adding 3D animation does not improve
the experience; it just adds another visual element competing with the wall. Content
relevance matters more than technical novelty for sustaining
engagement in voluntary, informal visits: Chen and Lai \cite{chenlai2021}, applying
the ARCS motivation framework to museum AR, found relevance to be the strongest
motivational driver among the dimensions they tested, ahead of novelty. That finding
argues for linking each point of interest to something the visitor already
recognises (a present-day street, a still-standing building) instead of relying on
technical spectacle to hold attention.

\subsection{Personalisation and Visitor Profiling}
\label{subsec:rw-personalisation}

The visitor models introduced in \S\ref{subsec:rw-visitor-models} were largely
descriptive: Falk's five identity-driven archetypes \cite{falk2009} and Pekarik,
Doering and Karns's four experience types \cite{pekarik1999} explain why visitors
differ, without by themselves prescribing what an interface should do about it.
Roussou and Katifori's work on the CHESS project \cite{roussoukatifori2018}
closes that gap directly in a cultural-heritage AR context: personas built from
visitor interest and prior knowledge were used to drive personalised interactive
storytelling at the Acropolis Museum and the Cité de l'espace, and the project
reports that matching narrative content to visitor type measurably increased
engagement relative to a one-size-fits-all script. A companion study from the same
research group \cite{antoniou2016} tested how to capture that visitor type in the
first place, comparing a short explicit quiz against indirect, behaviour-inferred
profiling at the Acropolis Museum, and is one of relatively few heritage AR studies
to treat the profiling mechanism itself as
a design problem worth evaluating.

The empirical case for personalising at all is, fortunately, no longer just
theoretical. He and Wang \cite{hewang2025} found that what they call
\emph{individuation} (the system's ability to match content to a visitor's expressed preferences) was the strongest single predictor of overall satisfaction
among all the stimuli they tested ($\beta = 0.286$, $p < 0.001$), ahead of
interactivity and authenticity. For a wall carrying upwards of a hundred points of
interest, this matters more than it would for a five-exhibit gallery: without some
mechanism for narrowing what is shown, a generalist visitor and a specialist visitor
receive an identical, and for one of them likely poorly calibrated, experience.

Following Antoniou et al.'s comparison \cite{antoniou2016}, this thesis adopts
explicit onboarding (a short, visitor-answered set of questions) over
inferred, behaviour-based profiling. The trade-off is real: explicit self-report
risks a visitor selecting a socially desirable answer instead of their genuine
interest level, where indirect profiling avoids that bias entirely. But explicit
onboarding offers three practical advantages that matter for a museum deployment:
it is fully transparent to the visitor, it produces deterministic behaviour with no
cold-start period before enough interaction data accumulates, and it requires no
behavioural data collection at all, which simplifies both privacy compliance and
museum partnership conversations considerably. Table~\ref{tab:profile-mapping} shows
how the resulting visitor categories map onto the identity and experience models
already discussed, and onto four concrete content-adaptation dimensions a heritage
wall application can vary.

\begin{table}[h]
\centering
\caption{Mapping established visitor models onto adjustable content dimensions for
a heritage-wall AR application.}
\label{tab:profile-mapping}
\small
\begin{tabular}{p{2.5cm}p{3.2cm}p{2.6cm}p{2.6cm}p{2.6cm}}
\toprule
\textbf{Visitor profile} & \textbf{Relates to} & \textbf{Vocabulary} & \textbf{Depth} & \textbf{Suggested route length} \\
\midrule
Tourist & Experience seeker \cite{falk2009}; object experience \cite{pekarik1999} & Accessible & Overview & Short, $\leq$4 POIs \\
Student & Explorer \cite{falk2009}; cognitive experience \cite{pekarik1999} & Accessible--intermediate & Moderate & Medium \\
Historian/specialist & Professional/hobbyist \cite{falk2009}; cognitive experience \cite{pekarik1999} & Academic & Specialist detail & Long, self-directed \\
Child/family & Facilitator \cite{falk2009}; social experience \cite{pekarik1999} & Simplified & Overview & Short, $\leq$2 POIs \\
\bottomrule
\end{tabular}
\end{table}

Whichever profile a visitor selects, the constraint established in
\S\ref{subsec:rw-ucd} still applies: personalisation should narrow depth and route
suggestions, not fragment the shared experience of a group visit where a tourist,
a student, and a child may be standing in front of the same building together. A
final, older but still relevant finding from Pekarik, Doering and Karns
\cite{pekarik1999} reinforces why asking matters in the first place: visitors
arrive already framing what they see through pre-existing motivations and interests;
they are not blank slates waiting for an institution's framing. Asking directly
(``what interests you most about this place's history?'') engages that
pre-existing motivation instead of overriding it with a generic introduction. That
is the underlying justification for treating onboarding as a substantive design
decision, not a perfunctory first screen.

















\subsection{Gamification, Storytelling and Temporal Visualisation}
\label{subsec:rw-gamification-storytelling}

Gamification (``the use of game design elements in non-game contexts,'' in Deterding et al.'s now-standard definition \cite{deterding2011}) is one of
the most frequently proposed engagement tools in heritage AR, but the evidence
for which specific elements actually work is more contested than design
guides often suggest. Souropetsis and Kyza's CompARe system \cite{souropetsis2025},
a scaffolded, QR-triggered AR learning environment built around a sixth-century
wall mosaic in Cyprus, shows that gamified AR \emph{can} produce real
learning gains when game elements are tightly coupled to real content
engagement; that qualifier matters, since a
broader line of research (Mekler et al.'s controlled study of individual game elements and Hanus and Fox's classroom comparison among them) finds
points, badges, and leaderboards by themselves to have an inconsistent, even
negative, effect on intrinsic motivation \cite{mekler2017,hanusfox2015}. The
practical implication for a wall carrying many points of interest is not
``add a leaderboard,'' but that any reward mechanic should stay contingent on
genuine content engagement, a point Liamruk et al.'s 2025 serious-game
evaluation in a cultural-tourism setting reinforces from a different angle,
framing gamified AR as most defensible when built around a learning
objective, not engagement metrics alone \cite{liamruk2025}.

Storytelling addresses a closely related but distinct problem: even where
content is informative in its own right, digital heritage tools tend to default to
information delivery instead of narrative, despite narrative being, on
Bruner's account, one of the most fundamental modes of human learning
\cite{gkatsou2018}. Xu's 2026 evaluation of the ``Hidden Stories'' AR
application at the Chinatown Storytelling Centre \cite{xu2026chinatown} is a
useful recent data point here: a marker-based AR experience combining
narrative content with gamified challenges was found to foster emotional
connection and improve knowledge retention relative to a non-narrative
baseline, in a real heritage institution, not a laboratory mock-up.
This connects back to the non-linear navigation principle established in
\S\ref{subsec:rw-ucd}: AR narrative naturally wants a fixed telling order,
while a physical wall is explored in whatever order a visitor happens to
walk it, and resolving that tension by allowing free entry into self-contained
thematic circuits, instead of forcing a single sequence, appears to be the
more visitor-respecting design choice on the evidence reviewed so far
\cite{gkatsou2018}.

A third engagement technique, temporal visualisation, layers one or more
historical states onto the present physical surface. The idea has real
precedent outside mobile AR: Berlin's Museum f\"ur Naturkunde uses
fixed-installation ``Jurascopes,'' look-through viewers that progressively
animate a dinosaur skeleton through muscle, skin, and motion, to make a
static fossil legible as a once-living animal~\cite{jurascope2026}. Printezis
and Koutsabasis's 2026 review of spatial AR storytelling explicitly
identifies ``time-based palimpsests'' (toggling between past and present states anchored to the same physical location) as a recurring narrative
pattern in this space \cite{printezis2026}, and the underlying logic
transfers reasonably well to mobile AR: a visitor-controlled toggle between
a wall's painted historical state and its present-day surroundings, not
an automatically-playing animation, lets each visitor set their own
pace. Sauter et al.'s AR-Explorer \cite{sauter2024}, which retrieves and
overlays historical photographs at their original outdoor locations, is the
closest mobile precedent for this kind of then-and-now layering, even though
its location-based anchoring differs from an image-anchored approach.
Whether such reconstructions should ultimately be rendered as 3D models
instead of flat image overlays is, on present evidence, not fully settled
for heritage AR specifically, but Kr\"uger, Palzer, and Bodemer's controlled
comparison of 3D and 2D AR visualisation for spatial learning found a clear
advantage for the 3D condition \cite{kruger2021}, which is a reasonable basis for
treating 3D reconstruction as a worthwhile, if resource-intensive, future
upgrade for a wall already carrying simpler image-based historical content.

Disaster reconstruction is a particularly sensitive special case of temporal
visualisation, and the 1755 Lisbon earthquake (which destroyed an estimated 85\% of the city's buildings, including many depicted in pre-1755 panoramic paintings) is the obvious narrative hook for any
heritage wall from this period. Simulating a mass-casualty historical event
in AR raises a design obligation beyond the purely technical: any
speculative or AI-assisted reconstruction content should be disclosed as
such, not presented as documented fact, and dramatic spectacle should
not be allowed to override historical accuracy or sensitivity to the human
cost of the event being depicted. Finally, connecting a historical scene to
its present-day surroundings (whether through a simple ``still exists / no longer exists'' marker or a fuller mapped overlay) has an
established conceptual precedent in the Museum of London's
\emph{Streetmuseum} app \cite{streetmuseum}, which overlaid archival photographs onto present-day
city locations via GPS-triggered outdoor AR. The distinction worth keeping in
mind, as in \S\ref{subsec:rw-large-scale}, is that Streetmuseum's
location-based outdoor anchoring does not transfer directly to an
image-anchored indoor wall; the conceptual pattern of bridging past and
present generalises across heritage walls of very different kinds, but the
underlying tracking mechanism has to be re-derived for each.

\subsection{Evaluation Methods for Heritage AR UX}
\label{subsec:rw-evaluation}

The studies reviewed across this chapter rely on a fairly small, recurring
toolkit of evaluation instruments. The System Usability Scale remains the
default usability measure across the field generally, often paired with a
domain-specific instrument such as the Museum Experience Scale used by Xu
\cite{xu2026chinatown}, or with a validated motivation framework such as the
ARCS model applied by Chen and Lai \cite{chenlai2021}. Comparative designs,
where an AR condition is evaluated directly against a traditional, non-AR
visit, remain comparatively rare; Bachiller, Monzo, and Rey's study
\cite{bachiller2023} is one of the few examples in this chapter to use that
structure instead of evaluating an AR prototype in isolation, and reads as
something closer to a methodological model than an
exception.

A second pattern is the trade-off between sample size
and ecological validity. Large, well-powered studies in this area, such as
He and Wang's 324-respondent survey \cite{hewang2025}, are generally
online or laboratory-based, not collected from a working museum
floor. Studies conducted in situ at a real heritage site, by contrast, tend
to run with considerably smaller samples: Souropetsis and Kyza's CompARe
evaluation and Xu's Chinatown Storytelling Centre study
\cite{souropetsis2025,xu2026chinatown} both work with samples in the
tens, not the hundreds, reflecting how much harder it is to recruit
and instrument real visitors than online participants. Neither approach is
strictly better; a survey of 324 respondents establishes a statistically
defensible effect size but says little about how a system behaves on an
actual gallery floor, while a 47-visitor in-situ study captures real
deployment conditions but cannot support strong statistical claims on its
own. Table~\ref{tab:evaluation-methods} summarises this trade-off across the
studies most directly relevant to this thesis.

\begin{table}[h]
\centering
\caption{Evaluation approaches used across the heritage AR literature
reviewed in this chapter.}
\label{tab:evaluation-methods}
\small
\begin{tabular}{p{2.6cm}p{2.6cm}p{1.6cm}p{2.2cm}p{3.4cm}}
\toprule
\textbf{Study} & \textbf{Instrument(s)} & \textbf{N} & \textbf{Setting} & \textbf{Design} \\
\midrule
He \& Wang \cite{hewang2025} & Path/SEM survey & 324 & Online & Cross-sectional \\
Bachiller et al.\ \cite{bachiller2023} & Satisfaction \& learning survey & --- & In-situ museum & AR vs.\ traditional comparison \\
Souropetsis \& Kyza \cite{souropetsis2025} & Custom + motivation measures & Small & In-situ heritage site & Single-condition, in-situ \\
Xu \cite{xu2026chinatown} & SUS + Museum Experience Scale & 47 & In-situ museum & Single-condition, in-situ \\
\bottomrule
\end{tabular}
\end{table}

This thesis sits closer to the in-situ end of that trade-off, on
purpose: the research gap identified in \S\ref{subsec:rw-gap}
specifically concerns the lack of evidence for large-scale heritage walls
evaluated with real visitors, not lab participants. Consistent with
the modest sample sizes typical of in-situ heritage AR studies generally,
this thesis's own evaluation (in-app analytics, brief exit surveys, and where feasible a short follow-up questionnaire) is reported as
exploratory and context-specific, not as a basis for strong
statistical generalisation. That is the framing the rest of this
sub-field's evaluation literature also tends to adopt.

















% ===================================================================
\section{Summary and Research Gap}
\label{sec:rw-summary}
% ===================================================================

\subsection{Comparative Overview of Related Work}
\label{subsec:rw-comparative}

Read across the three preceding sections, not within each one, and a
fairly consistent shape emerges: a great deal is known about what visitors want
from a heritage encounter (\S\ref{sec:rw-heritage}), a great deal is known
about how to track and render AR content on a heritage object
(\S\ref{sec:rw-ar-ch}), and a great deal is known about how to keep that
content engaging without overwhelming the visitor
(\S\ref{sec:rw-ux}). What is comparatively rare is a single system, or even
a single published study, that combines all three at once on a surface
large enough to make them all matter simultaneously. Most of the strongest
evidence reviewed in this chapter comes from studies built around one
object, one room, or one façade: Recupero et al.'s Ara Pacis tour
\cite{recupero2019}, Gkatsou's Alf\^{a}ndega Museum thesis
\cite{gkatsou2018}, Bachiller et al.'s technological heritage museum
comparison \cite{bachiller2023}, Souropetsis and Kyza's wall-mosaic
deployment \cite{souropetsis2025}, and Xu's Chinatown Storytelling Centre
study \cite{xu2026chinatown} are all, in this sense, single-site, modest-scale
systems, however methodologically strong individually. Bauer, Hagler, and
Kocur's mural application \cite{bauer2025} is the closest in raw surface
scale, but its outdoor GPS anchoring is unavailable indoors and its
evaluation does not address personalisation at all. Table
\ref{tab:comparative-overview} makes this pattern concrete across six
representative works and the design this thesis pursues.

\begin{table}[h]
\centering
\caption{Representative related work compared against the dimensions this
thesis combines. \checkmark{} = addressed, $\sim$ = partially addressed,
--- = not addressed.}
\label{tab:comparative-overview}
\small
\begin{tabular}{p{3.8cm}p{1.6cm}p{1.6cm}p{1.6cm}p{1.6cm}p{1.8cm}}
\toprule
\textbf{Work} & \textbf{Large, multi-POI scale} & \textbf{Portable (no fixed install)} & \textbf{Persona-lisation} & \textbf{Non-linear by design} & \textbf{In-situ evaluation} \\
\midrule
Recupero et al.\ \cite{recupero2019} & --- & \checkmark & --- & $\sim$ & \checkmark \\
Gkatsou \cite{gkatsou2018} & --- & \checkmark & --- & \checkmark & $\sim$ \\
Bauer et al.\ \cite{bauer2025} & \checkmark & \checkmark & --- & $\sim$ & \checkmark \\
Souropetsis \& Kyza \cite{souropetsis2025} & --- & \checkmark & --- & $\sim$ & \checkmark \\
Xu \cite{xu2026chinatown} & --- & \checkmark & --- & $\sim$ & \checkmark \\
\textbf{This thesis} & \checkmark & \checkmark & \checkmark & \checkmark & \checkmark \\
\bottomrule
\end{tabular}
\end{table}

No row above checks every column, and that gap is not evenly distributed
across the dimensions either. Scale and portability are each individually
solved (Bauer et al.\ show a large surface tracked without a fixed installation, just outdoors, not indoors) but personalisation is the
column where the literature is weakest, exactly where it is most needed:
none of the in-situ heritage AR deployments reviewed in this chapter
adapts content depth or route suggestions to who the visitor is, even
though \S\ref{subsec:rw-personalisation} showed personalisation to be among
the strongest predictors of visitor satisfaction once it is actually tested
\cite{hewang2025}. Evaluation, meanwhile, is where ambition and feasibility
trade off most visibly: the studies with the largest, most statistically
defensible samples \cite{hewang2025} are not conducted on a real museum
floor, and the studies that are tend to run with the tens-of-participants
samples discussed in \S\ref{subsec:rw-evaluation}, not hundreds.

\subsection{Positioning of This Thesis}
\label{subsec:rw-positioning}

This thesis is positioned at the intersection that Table
\ref{tab:comparative-overview} shows to be largely empty: a mobile,
portable AR system, validated across the entire length of a large,
multi-POI heritage wall, not a single object or room; one that adapts
content to the visitor instead of presenting one fixed script; and one that
is evaluated, however modestly, with real visitors in a real heritage
setting, not only in a lab. None of the individual ingredients is
new (mobile AR tracking, visitor personalisation, and in-situ evaluation are each well studied in isolation, as this chapter has shown at length)
but their combination, applied to a continuous wall at the scale of tens
of metres and on the order of a hundred or more points of interest, does
not appear to have a close published precedent.

The Grande Panorama de Lisboa remains this thesis's central motivating
case, and its roughly 150 points of interest are the concrete number
behind every design decision discussed above. But the framework this
thesis develops is pitched, on purpose, at the level of ``a large
heritage wall,'' not ``the Grande Panorama specifically.'' Chafariz Velho,
used extensively during development while the Museu Nacional do Azulejo
was closed for renovation, served as a harder generalisation test: outdoor, not indoor; circular,
not U-shaped; and therefore demanding yaw tracking and variable natural lighting that the Panorama
itself never requires (\S\ref{subsec:rw-ar-concepts-tracking}). A framework
validated only on one flat, indoor, climate-controlled wall would say very
little about whether it generalises; validating central design decisions
against both surfaces is what allows this thesis to claim, with reasonable
confidence, that its contribution extends beyond a single building's
collection.

Concretely, the literature reviewed in this chapter motivates four design
commitments carried forward into Chapter~3. First, tracking is approached
as a multi-zone, image-anchored problem, not a single-marker or
single-scan one, following directly from the scale and acquisition
constraints established in \S\ref{subsec:rw-large-scale}. Second, content
is structured as short, non-linear, visitor-initiated circuits instead of
one mandatory sequence, following the satiation and non-linear-navigation
evidence in \S\ref{subsec:rw-visitor-models} and
\S\ref{subsec:rw-ucd}. Third, a short explicit onboarding step adapts
depth and route suggestions to a small set of visitor profiles, following
\S\ref{subsec:rw-personalisation}'s evidence that personalisation is
disproportionately under-served relative to its measured impact on
satisfaction. Fourth, evaluation is planned as an in-situ,
modest-sample exercise (in-app analytics, brief exit surveys, and where feasible a short follow-up questionnaire) in keeping with
\S\ref{subsec:rw-evaluation}'s observation that this is how most of the
field's strongest in-situ work is actually conducted. It does not
attempt to claim a statistical power this kind of study cannot
realistically support. Chapter~3 develops each of these four commitments
into a concrete system design.